%----------------------------------------------------------------------------------------
%	CAPÍTULO 1
%	Este capítulo trae ejemplos de todo lo que necesitas en APA (citas, tablas,
%	figuras, citas largas, recuadros). Úsalo como guía y reemplaza el contenido.
%----------------------------------------------------------------------------------------

\chapter{Generalidades}

Texto introductorio del capítulo. En APA cada párrafo empieza con sangría de 1,27 cm; para un párrafo nuevo basta con dejar una línea en blanco en el código. \lipsum[1]

%------------------------------------------------

\section{Cómo citar en APA}

Las citas se generan solas a partir de \texttt{referencias.bib}:

\begin{itemize}
\item Cita narrativa: \textcite{apa2020} establece las normas de esta edición.
\item Cita parentética: las normas se actualizaron en su séptima edición \parencite{apa2020}.
\item Con número de página: \parencite[p.~45]{apellido2023}.
\item Varias fuentes a la vez: \parencite{apellido2023,autor2024}.
\item Tres o más autores (sale \enquote{et al.} automáticamente): \textcite{autor2024}.
\item Institución como autor: \parencite{dane2024}.
\item Otros tipos de fuente (capítulo, tesis, web, ley): \parencite{capitulo2022,tesis2021,web2025,ley617}.
\end{itemize}

Para una cita textual corta (menos de 40 palabras) usa comillas: según \textcite[p.~12]{apellido2023}, \enquote{la cita textual va entre comillas y con página}.

Las citas de 40 palabras o más van en bloque, con sangría y sin comillas:

\begin{citalarga}
Este es un ejemplo de cita en bloque. APA indica que las citas textuales de 40 palabras o más se presentan en un párrafo independiente, con sangría de 1,27 cm desde el margen izquierdo, sin comillas, y con la cita entre paréntesis al final, después del punto. \parencite[p.~171]{apa2020}
\end{citalarga}

%------------------------------------------------

\section{Tablas en formato APA}

Las tablas se mencionan en el texto antes de aparecer, por ejemplo: la Tabla~\ref{tab:ejemplo} resume los datos. El número va en negrita, el título en cursiva, ambos encima de la tabla; la nota va debajo.

\begin{table}[htbp]
\caption{Título de la tabla en cursiva y con mayúscula inicial}
\label{tab:ejemplo}
\begin{tabular*}{\linewidth}{@{\extracolsep{\fill}} l c c}
\toprule
\textbf{Municipio} & \textbf{Indicador 1} & \textbf{Indicador 2}\\
\midrule
Municipio A & 0,32 & 56,2 \\
Municipio B & 0,15 & 91,0 \\
Municipio C & 0,09 & 29,6 \\
\bottomrule
\end{tabular*}
\nota{Elaboración propia con datos de \textcite{dane2024}.}
\end{table}

%------------------------------------------------

\section{Figuras en formato APA}

Las figuras (gráficos, imágenes, diagramas, capturas de pantalla) siguen el mismo formato, como en la Figura~\ref{fig:ejemplo}.

\begin{figure}[htbp]
\caption{Título de la figura en cursiva}
\label{fig:ejemplo}
\centering
\includegraphics[width=0.6\linewidth]{placeholder}
\nota{Tomado de \textcite{autor2024}. Si la figura es propia, escribe \enquote{Elaboración propia}.}
\end{figure}

%------------------------------------------------

\section{Subtítulos y listas}

\subsection{Subtítulo de nivel 3}

Los capítulos equivalen al nivel 1 de APA, las secciones (\verb|\section|) al nivel 2 y las subsecciones (\verb|\subsection|) al nivel 3.

\nivelcuatro{Encabezado de nivel 4} Si necesitas un nivel más, usa \verb|\nivelcuatro{...}|: va en negrita, con sangría y en la misma línea del párrafo.

\nivelcinco{Encabezado de nivel 5} Y \verb|\nivelcinco{...}| va en negrita cursiva.

\subsection{Listas}

\begin{enumerate}
\item Primer elemento.
\item Segundo elemento.
\item Tercer elemento.
\end{enumerate}

%------------------------------------------------

\section{Recuadros de la plantilla}

La plantilla conserva los recuadros de colores del diseño original, ahora en español.

\begin{definition}[Gestión pública]
Escribe aquí la definición de un concepto clave del libro \parencite{apellido2023}.
\end{definition}

\begin{example}[Nombre del ejemplo]
Texto de un ejemplo práctico.
\end{example}

\begin{exercise}
Una pregunta o actividad para el lector.
\end{exercise}

\begin{remark}
Una nota o aclaración al margen del texto principal.
\end{remark}
